\section{The split branch on the tetraquadric}\label{sec:splitbranch}
The ledger of the next section records two failures at the fixed datum
that do not depend on the complex structure: the residual Abelian factor
of Section~\ref{sec:extraU1} and the Yukawa deficiency of
Proposition~\ref{prop:rankfailure}. This section asks whether another
choice of five line bundles on the same manifold could avoid them while
staying at the split locus. The first failure cannot be avoided by any
sum of line bundles on any Calabi--Yau threefold with at most four K\"ahler
parameters; the second is a property of the charge matrix, and the
selection rules decide it in finitely many steps for each datum.

\subsection{The residual Abelian factor is forced}\label{sec:forced}
Let $X$ be a Calabi--Yau threefold, $h=h^{1,1}(X)$, and let $J_1,\dots,J_h$
be divisor classes spanning $H^{1,1}(X,\Q)$. A sum of line bundles
$V=\bigoplus_{a=1}^nL_a$ with $c_1(L_a)=\sum_ik_a^iJ_i$ has the integer
charge matrix $K=(k_a^i)\in\Z^{h\times n}$, whose columns are the vectors
$k_a$. For a K\"ahler class $J$ write $\kappa_i(J)=\int_XJ_iJ^2$.

\begin{theorem}[The residual factor]\label{thm:forced}
Suppose $c_1(V)=0$ and $V$ is slope-polystable with respect to $J$. Then
\begin{enumerate}[label=\textup{(\roman*)},leftmargin=*,itemsep=1pt]
\item $\rk K\le h-1$;
\item the traceless kernel $\{q\in\Q^n:\ Kq=0,\ \sum_aq_a=0\}$ has dimension
at least $n-h$;
\item at the split locus, with vanishing expectation values of all charged
singlets, exactly $n-1-\rk K\ge n-h$ of the $n-1$ traceless Abelian gauge
bosons of the structure group $S(U(1)^n)$ are massless.
\end{enumerate}
In particular, for $n=5$ and $h\le4$ at least one Abelian gauge boson
beyond the Standard Model stays massless at every split point of every
rank-five sum of line bundles on $X$, and condition C3, read as the
visible gauge group being exactly the Standard Model group, fails there.
\end{theorem}
\begin{proof}
The slope of $L_a$ is $\mu_J(L_a)=\int_Xc_1(L_a)J^2=\sum_ik_a^i\kappa_i$.
A polystable sum has all summands of slope $\mu_J(V)$, and
$\mu_J(V)=\int_Xc_1(V)J^2=0$. Thus $K^{T}\kappa=0$. The vector $\kappa$ is
nonzero, because $\sum_it_i\kappa_i=\int_XJ^3>0$ when $J=\sum_it_iJ_i$.
Hence the rows of $K$ satisfy one nontrivial linear relation and
$\rk K\le h-1$, which is (i). Since $c_1(V)=\sum_ac_1(L_a)=0$, the vector
$\one=(1,\dots,1)$ lies in $\ker K$, which has dimension
$n-\rk K\ge n-h+1$; the hyperplane $\sum_aq_a=0$ does not contain $\one$,
so it meets $\ker K$ in a space of dimension at least $n-h$, which is (ii).

For (iii) we use the Green--Schwarz masses of the Abelian gauge bosons as
computed by Anderson, Gray, Lukas and Palti \cite[Section~2.7]{AGLP12}.
The axions of the $h$ K\"ahler moduli and the universal axion shift under
the $U(1)$ gauge transformations with coefficients $k_a^i$ and
$\beta_ik_a^i$ respectively, and the mass matrix of the gauge bosons
$A_a$ in the basis of the $n$ summands is
$M_{ab}=k_a^{T}\widetilde Gk_b$, where
$\widetilde G=G+\beta\beta^{T}/(4s^2)$, $G$ is the positive definite
metric on the K\"ahler moduli and $s$ the dilaton. The matrix
$\widetilde G$ is positive definite, so the quadratic form
$q\mapsto q^{T}Mq=(Kq)^{T}\widetilde G(Kq)$ vanishes exactly on $\ker K$:
the combination $\sum_aq_aA_a$ is massless if and only if $Kq=0$. On the
traceless vectors, which parametrize the gauge bosons of $S(U(1)^n)$, the
massless combinations form the space of (ii), of dimension
$n-1-\rk K$. This is the count of \cite[eq.~(2.31)]{AGLP12}, and (i) gives
the lower bound $n-h$, which is the statement of \cite[eq.~(2.32)]{AGLP12}
that fewer than five K\"ahler moduli leave at least one $U(1)$ massless.
A nonzero traceless $q$ in $\ker K$ acts nontrivially on some summand, so
the surviving gauge boson is not a redundant generator.
\end{proof}

\begin{remark}[Reading the fixed datum]\label{rem:forcedfixed}
For the matrix \eqref{eq:bundle}, $\rk K=3=h-1$, the largest value
Theorem~\ref{thm:forced} allows, and $\ker K$ is spanned by $\one$ and the
vector $q$ of \eqref{eq:extrau1}. The residual factor of
Section~\ref{sec:extraU1} is therefore not a defect of the choice
\eqref{eq:bundle}; it is the minimum that the tetraquadric, with
$h^{1,1}=4$, imposes on every rank-five split datum. The polarization
locus is tied to the same number. The constraint $K^{T}\kappa=0$ confines
$\kappa$ to a space of dimension $h-\rk K=1$, the line through
$(1,1,1,1)$, and for $J=\sum_it_iJ_i$ with all $t_i>0$ the intersection
numbers of Proposition~\ref{prop:cover} give
$\kappa_i=4\,e_2(t_j:j\ne i)=4\bigl(e_2(t)-t_i(e_1(t)-t_i)\bigr)$, where
$e_1,e_2$ are the elementary symmetric functions. All four $\kappa_i$ are
equal exactly when $t_i(e_1(t)-t_i)$ takes one value, so each $t_i$ is a
root of $x(e_1-x)=c$, hence $t_i\in\{c',e_1-c'\}$ for one $c'$; if $m$ of
the coordinates equal $c'$ and $4-m$ equal $e_1-c'$, summing gives
$(m-3)e_1=(2m-4)c'$, which has a solution with all $t_i>0$ only for $m=0$
or $m=4$, that is, all coordinates equal. Thus the ray $J=\lambda\sum_iJ_i$
of Section~\ref{sec:pointfailure} is the whole locus of K\"ahler classes
at which \eqref{eq:bundle} is polystable, and the flat classical direction
of C7 and the massless generator of C3 are two readings of the one
equality $4-\rk K=1$. For $n=h+1$ the two counts $n-1-\rk K$ and
$h-\rk K$ agree for every split datum.
\end{remark}

\subsection{Selection rules on eleven split data}\label{sec:eleven}
Let $X\subset(\PP^1)^4$ be any smooth tetraquadric and $V=\bigoplus_{a=1}^5
\cO_X(k_a)$ with $c_1(V)=0$. Under the structure group $S(U(1)^5)$ the
charged fields carry the charges of Section~\ref{sec:extraU1}, in the
coordinates $e_1,\dots,e_5$ modulo the diagonal $\one$, as listed in
Table~\ref{tab:charges}.
\begin{table}[t]\small\centering
\caption{The charged fields of a rank-five split datum under
$S(U(1)^5)$.}\label{tab:charges}
\begin{tabular}{@{}llll@{}}
\toprule
field & cohomology & charge & representation\\\midrule
family $\mathbf{10}$ & $H^1(L_a)$ & $e_a$ & $\mathbf{10}$ of $SU(5)$\\
family or Higgs $\overline{\mathbf 5}$ & $H^1(L_aL_b)$, $a<b$ & $e_a+e_b$ & $\overline{\mathbf5}$\\
Higgs $\mathbf 5$ & $H^2(L_aL_b)$, $a<b$ & $-e_a-e_b$ & $\mathbf 5$\\
singlet & $H^1(L_aL_b^{-1})$, $a\ne b$ & $e_a-e_b$ & $\mathbf 1$\\
\bottomrule
\end{tabular}
\end{table}
A term of the perturbative superpotential is a product of charged fields
times a monomial in the singlets, and it is gauge invariant exactly when
its total charge is a multiple of $\one$. For the down-type coupling
$\mathbf{10}_a\,\overline{\mathbf5}_{bc}\,\overline{\mathbf5}_{de}$ the charge
$Q=e_a+e_b+e_c+e_d+e_e$ has coordinate sum five, so invariance after
insertion of singlets requires $\one-Q$ to be a nonnegative integer
combination of the available singlet charges; for the up-type coupling
$\mathbf{10}_a\mathbf{10}_b\mathbf 5_{de}$ the same is required of
$e_d+e_e-e_a-e_b$. Which groups are nonzero depends on $X$, but on every
smooth tetraquadric the restriction sequence
$0\to\cO_A(k-2\one)\to\cO_A(k)\to\cO_X(k)\to0$ gives
\begin{equation}\label{eq:koszulbound}
 h^i(X,\cO_X(k))\le h^i(A,\cO_A(k))+h^{i+1}(A,\cO_A(k-2\one)),
\end{equation}
with the ambient groups given by the K\"unneth formula on $A=(\PP^1)^4$.
We call a field \emph{possibly present} when the bound
\eqref{eq:koszulbound} for its group is positive. A coupling forbidden
with every possibly present singlet is forbidden on every smooth
tetraquadric in the class; a coupling allowed by the charges is only not
excluded by them.

\begin{lemma}[Flow criterion]\label{lem:flow}
Let $E$ be a set of ordered pairs $(p,q)$ of distinct elements of
$N=\{1,\dots,n\}$ and $v\in\Z^n$ with $\sum_pv_p=0$. The following are
equivalent.
\begin{enumerate}[label=\textup{(\alph*)},leftmargin=*,itemsep=1pt]
\item $v=\sum_{(p,q)\in E}m_{pq}(e_p-e_q)$ with integers $m_{pq}\ge0$.
\item For every $B\subset N$ that no pair of $E$ leaves, that is, $p\in B$
and $(p,q)\in E$ imply $q\in B$, one has $\sum_{p\in B}v_p\le0$.
\end{enumerate}
\end{lemma}
\begin{proof}
A pair $(p,q)$ contributes $e_p-e_q$, whose sum over $B$ is $1$ if
$p\in B\not\ni q$, $-1$ if $q\in B\not\ni p$ and $0$ otherwise. When no pair
leaves $B$ the first case does not occur, so (a) gives
$\sum_Bv\le0$, which is (b).

Conversely, form the network with a source $s$, a sink $t$, the vertices
of $N$, arcs $s\to p$ of capacity $v_p$ for $v_p>0$, arcs $q\to t$ of
capacity $-v_q$ for $v_q<0$, and arcs $p\to q$ of unbounded capacity for
$(p,q)\in E$. Let $S\ni s$, $T\ni t$ be a cut of finite capacity and
$B=S\cap N$. No pair of $E$ leaves $B$, since such an arc would cross the
cut with unbounded capacity. The capacity of the cut is
\[
 \sum_{p\notin B,\ v_p>0}v_p+\sum_{q\in B,\ v_q<0}(-v_q)
 =\sum_{v_p>0}v_p-\sum_{p\in B}v_p\ \ge\ \sum_{v_p>0}v_p
\]
by (b). The cut $S=\{s\}$ has capacity $\sum_{v_p>0}v_p$, so this is the
minimum, and the maximal flow theorem of Ford and Fulkerson \cite{FF56}
gives an integral flow of that value, because all capacities are
integers. It saturates every arc at $s$ and at $t$. Let $m_{pq}$ be the
flow on the arc $p\to q$. Conservation at $p\in N$ says that
$\sum_qm_{pq}-\sum_qm_{qp}$ equals the inflow from $s$ minus the outflow
to $t$, which is $v_p$. This is (a).
\end{proof}

The criterion is decided by inspecting the $2^5$ subsets; the cases below
were also checked by a maximal flow computation.\footnote{The program
\texttt{code/yukawa\_charges.py} of the repository named in the note under
the abstract lists, for each datum, the possibly present fields by
\eqref{eq:koszulbound}, decides the two couplings by a maximal flow and
prints either a closed set $B$ violating (b) or an explicit decomposition
(a). Its output is the file \texttt{yukawa\_charges\_results.json}.}

\begin{proposition}[Eleven split data]\label{prop:eleven}
Consider the matrix \eqref{eq:bundle} and the ten rank-five sums of line
bundles with entries in $[-2,2]$ listed in Table~\ref{tab:eleven}, which are
the classes, up to permutation of the four factors of $(\PP^1)^4$, found
by the scan of \cite{Anaxia} to satisfy $c_1(V)=0$, the anomaly inequality,
the index $-3|G|$ for a free group $G$ of the stated order, a K\"ahler
class with all slopes zero, and the upstairs spectrum tests of
Section~\ref{sec:fixedbundle}. With all possibly present singlets
inserted to every order:
\begin{enumerate}[label=\textup{(\roman*)},leftmargin=*,itemsep=1pt]
\item for \eqref{eq:bundle} every down-type coupling vanishes on every
smooth tetraquadric, while the up-type couplings
$\mathbf{10}_2\mathbf{10}_5\mathbf 5_{25}$ and, with one singlet of charge
$e_5-e_2$, $\mathbf{10}_2\mathbf{10}_2\mathbf 5_{25}$ and
$\mathbf{10}_5\mathbf{10}_5\mathbf 5_{25}$ are allowed by the charges;
\item for the classes $9$ and $10$ every up-type and every down-type
coupling vanishes; for the classes $3$, $4$ and $5$ every up-type coupling
vanishes;
\item for the classes $1$, $2$, $6$, $7$ and $8$ both couplings are
allowed by the charges, with the singlet insertions of
Table~\ref{tab:eleven}.
\end{enumerate}
Hence six of the eleven split data fail C8 to all orders of the
perturbative singlet expansion on every smooth tetraquadric, and the five
others are not excluded by the selection rules.
\end{proposition}

\begin{table}[t]\small\centering
\caption{The eleven split data. Rows of $K$ are the vectors $k_a$; the
$\mathbf{10}$ column lists the summands $L_a$ with $H^1(L_a)$ possibly
present; ``down'' and ``up'' give a closed set $B$ of Lemma~\ref{lem:flow}
forbidding every coupling of that type, or one allowed coupling with the
singlet insertions it needs ($p\to q$ denotes a singlet of charge
$e_p-e_q$).}\label{tab:eleven}
\begin{adjustbox}{max width=\linewidth}
\begin{tabular}{@{}clcll@{}}
\toprule
datum & rows $k_a$ & $\mathbf{10}$ & down & up\\\midrule
\eqref{eq:bundle} & $(-1,0,0,1),(-1,-3,2,2),(0,1,-1,0),(1,1,-1,-1),(1,1,0,-2)$ & $L_2,L_5$ & $B=\{1\}$, $B=\{3\}$ & $\mathbf{10}_2\mathbf{10}_5\mathbf5_{25}$\\
1 ($|G|=2$) & $(-2,1,1,-1),(-1,-1,1,1),(-1,-1,1,1),(2,0,-2,1),(2,1,-1,-2)$ & $L_4$ & $\mathbf{10}_4\overline{\mathbf5}_{12}\overline{\mathbf5}_{12}$; $3\to1$, $5\to2$ & $\mathbf{10}_4\mathbf{10}_4\mathbf5_{12}$; $2\times(5\to4)$, $1\to5$, $2\to5$\\
2 ($|G|=4$) & $(-2,1,0,-1),(-1,-1,1,1),(-1,-1,1,1),(2,0,-2,1),(2,1,0,-2)$ & $L_4,L_5$ & $\mathbf{10}_4\overline{\mathbf5}_{12}\overline{\mathbf5}_{12}$; $3\to1$, $5\to2$ & $\mathbf{10}_4\mathbf{10}_4\mathbf5_{15}$; $2\times(5\to4)$, $1\to5$\\
3 ($|G|=2$) & $(-1,0,0,1),(-1,0,0,1),(-1,0,0,1),(1,-2,1,-1),(2,2,-1,-2)$ & $L_5$ & $\mathbf{10}_5\overline{\mathbf5}_{12}\overline{\mathbf5}_{34}$; none & $B=\{4\}$\\
4 ($|G|=2$) & $(-1,-1,1,0),(-1,0,0,1),(-1,0,0,1),(1,-1,0,0),(2,2,-1,-2)$ & $L_5$ & $\mathbf{10}_5\overline{\mathbf5}_{14}\overline{\mathbf5}_{23}$; none & $B=\{2,4\}$\\
5 ($|G|=2$) & $(-1,0,1,0),(-1,0,1,0),(0,-2,0,1),(1,1,-2,0),(1,1,0,-1)$ & $L_3,L_4$ & $\mathbf{10}_3\overline{\mathbf5}_{12}\overline{\mathbf5}_{45}$; none & $B=\{1\}$\\
6 ($|G|=4$) & $(-1,0,1,0),(-1,1,0,0),(0,-2,1,1),(1,0,-2,1),(1,1,0,-2)$ & $L_3,L_4,L_5$ & $\mathbf{10}_3\overline{\mathbf5}_{12}\overline{\mathbf5}_{12}$; $4\to1$, $5\to2$ & $\mathbf{10}_3\mathbf{10}_3\mathbf5_{25}$; $2\times(4\to3)$, $2\times(5\to4)$, $2\to5$\\
7 ($|G|=4$) & $(-1,0,0,1),(-1,0,0,1),(0,-2,1,1),(1,1,-2,-1),(1,1,1,-2)$ & $L_3,L_5$ & $\mathbf{10}_3\overline{\mathbf5}_{12}\overline{\mathbf5}_{12}$; $5\to1$, $4\to5$, $5\to2$ & $\mathbf{10}_3\mathbf{10}_3\mathbf5_{12}$; $2\times(5\to3)$, $2\times(4\to5)$, $1\to4$, $2\to4$\\
8 ($|G|=4$) & $(-1,-1,1,1),(-1,1,0,0),(0,-2,1,1),(1,1,-2,0),(1,1,0,-2)$ & $L_3,L_4,L_5$ & $\mathbf{10}_3\overline{\mathbf5}_{12}\overline{\mathbf5}_{12}$; $4\to1$, $5\to2$ & $\mathbf{10}_4\mathbf{10}_4\mathbf5_{24}$; $2\to4$\\
9 ($|G|=4$) & $(-1,0,1,0),(-1,1,-2,0),(0,-1,0,1),(0,-1,0,1),(2,1,1,-2)$ & $L_5$ & $B=\{1\}$ & $B=\{1,2,3\}$\\
10 ($|G|=4$) & $(-1,0,1,0),(-1,1,-2,0),(0,-1,0,1),(1,-1,-1,1),(1,1,2,-2)$ & $L_5$ & $B=\{1\}$ & $B=\{1,2\}$\\
\bottomrule
\end{tabular}
\end{adjustbox}
\end{table}

\begin{proof}
For each datum the possibly present fields follow from
\eqref{eq:koszulbound} and the K\"unneth formula
$h^0(\PP^1,\cO(m))=m+1$ for $m\ge0$, $h^1(\PP^1,\cO(m))=-m-1$ for $m\le-2$,
all other groups zero. For \eqref{eq:bundle} the possibly present
$\mathbf{10}$'s lie in $H^1(L_2)$ and $H^1(L_5)$, the
$\overline{\mathbf5}$'s in $H^1$ of $L_1L_2$, $L_2L_4$, $L_2L_5$, $L_4L_5$,
the $\mathbf5$'s in $H^2$ of $L_1L_2$, $L_2L_5$, and the singlets have
charges $e_p-e_q$ with $p\in\{2,4,5\}$ and $(p,q)$ one of
$(2,1),(2,3),(2,4),(2,5),(4,2),(5,1),(5,2),(5,3)$. No singlet has $p=1$ or
$p=3$, so the sets $B=\{1\}$ and $B=\{3\}$ are left by no pair. For a
down-type coupling $v=\one-Q$ has $v_3=1-Q_3=1$, since no charged field
carries the coordinate $e_3$; Lemma~\ref{lem:flow}(b) fails at $B=\{3\}$
and the coupling is forbidden. This recovers
Proposition~\ref{prop:singletobstruction} for every smooth tetraquadric
with the matrix \eqref{eq:bundle}. For the up-type couplings
$Q=e_a+e_b-e_2-e_5$ with $a,b\in\{2,5\}$; $Q=0$ for $\{a,b\}=\{2,5\}$, and
for $a=b=2$ the singlet $(5,2)$ gives $e_5-e_2=-Q$, for $a=b=5$ the singlet
$(2,5)$. This is (i).

For (ii) and (iii) the same two steps are carried out for each row of
Table~\ref{tab:eleven}. Where the table gives a set $B$, the reader
checks that no possibly present singlet $(p,q)$ has $p\in B$, $q\notin B$,
and that $\sum_Bv>0$ for every choice of the charged fields: for the
classes $9$ and $10$ the summand $L_1$ carries no $\mathbf{10}$ and lies in
no possibly present $\overline{\mathbf5}$, so $v_1=1$ for every down-type
coupling, while for the up-type couplings $B$ contains the coordinates of
every possibly present $\mathbf 5$ and no $\mathbf{10}$, so
$\sum_Bv=2>0$; for the classes $3$, $4$, $5$ the set $B$ contains a
coordinate of every possibly present $\mathbf5$ and misses the
coordinates of the $\mathbf{10}$'s. Where the table gives a coupling with
insertions, the listed singlets are possibly present and their charges
add to $\one-Q$, respectively $-Q$, which is (a) of the lemma. For the
classes $3$, $4$ and $5$ the down-type coupling listed has $Q=\one$ and
needs no insertion.
\end{proof}

\begin{remark}[What the proposition does and does not say]\label{rem:elevenscope}
The proposition decides a necessary condition. For the six excluded data
the conclusion is unconditional within the perturbative singlet
expansion: no complex structure, no polarization and no singlet
expectation values restore the missing matrices, since the bound
\eqref{eq:koszulbound} holds on every smooth member and
Lemma~\ref{lem:flow}(b) is violated. For the five remaining data nothing
is asserted beyond the absence of this obstruction. Their down-type
couplings arise only at second or third order in the singlets, the
singlet groups may vanish on a given member, the free action and its
projection of the spectrum have not been computed, and the stability of
the bundle after the singlets acquire expectation values, which is the
same extension problem as in Section~\ref{sec:pointfailure}, is open.
By Theorem~\ref{thm:forced} every one of the eleven data keeps a massless
Abelian gauge boson at its split point, so the five classes are candidates
for the extension branch, not for the split locus. The next subsection
decides, for each of the eleven data, whether the smallest non-Abelian
deformations of the split bundle remove that gauge boson.
\end{remark}

\subsection{The extension branch}\label{sec:extbranch}
Throughout this subsection $X\subset(\PP^1)^4$ is a smooth tetraquadric,
$V=\bigoplus_{a=1}^5L_a$ is one of the eleven data of
Table~\ref{tab:eleven} other than the class $3$, so that $\rk K=3$, and
$u\in\Z^5$ is a generator of the traceless kernel
$\{q:Kq=0,\ \sum_aq_a=0\}$, which has dimension one by
Theorem~\ref{thm:forced}(ii); the vectors $u$ are listed in
Table~\ref{tab:twoblocks}. For a K\"ahler class $J=\sum_it_iJ_i$,
$t_i>0$, we write $\mu_a(J)=\mu_J(L_a)=\sum_ik_a^i\kappa_i(J)$ and
$\mu(J)=(\mu_1,\dots,\mu_5)=K\kappa(J)$. The vector $\mu(J)$ lies in the
column space of $K$, which has dimension three and is contained in the
hyperplane $\sum_a\mu_a=0$.

The bundles that replace $V$ when the charged singlets of
Section~\ref{sec:pointfailure} acquire expectation values are the
polystable bundles $E=\bigoplus_jE_j$ of slope zero whose stable summands
$E_j$ are iterated extensions of the line bundles $L_a$, $a\in B_j$, for a
partition $\{B_j\}$ of $\{1,\dots,5\}$; they have the Chern character of
$V$, and we call the partition the \emph{shape} of $E$. The split bundle
has the shape with five blocks. The dictionary between such bundles and the
four-dimensional theory is the one of the stability walls of Anderson,
Gray, Lukas and Ovrut \cite{AGLO09,AGL09}: an extension class in
$H^1(L_pL_q^{-1})=\operatorname{Ext}^1(L_q,L_p)$ is the expectation value
of a singlet of charge $e_p-e_q$, the slopes of the summands are the
D-terms of the Abelian factors, and a non-split extension exists as a
polystable bundle exactly where the subbundle has the smaller slope.

\begin{lemma}[The Abelian factor on a branch]\label{lem:branchfactor}
Let $E=\bigoplus_{j=1}^mE_j$ be polystable of slope zero with the shape
$\{B_j\}$. The structure group of $E$ reduces to
$S\bigl(\prod_jU(|B_j|)\bigr)$, whose centre is parametrized by the
traceless vectors $q\in\Q^5$ constant on every block. The Abelian gauge
boson of such a generator $q$ is massless at $E$ if and only if $Kq=0$.
Consequently $E$ keeps a massless Abelian gauge boson beyond the
$SU(5)$ factor exactly when $u$ is constant on every block of its shape.
\end{lemma}
\begin{proof}
Each $E_j$ is a holomorphic bundle of rank $|B_j|$ with a Hermitian
Yang--Mills connection, so the connection on $E$ has holonomy in
$\prod_jU(|B_j|)$, and in the subgroup with determinant one because
$c_1(E)=c_1(V)=0$ and the connection is traceless. The centre of this group
consists of the diagonal matrices with one scalar per block, that is the
vectors $q$ constant on the blocks, traceless for the determinant
condition. The Green--Schwarz masses of these generators are computed as in
the proof of Theorem~\ref{thm:forced}: the axions shift under the generator
$q$ by the first Chern class $\sum_jq_{B_j}c_1(E_j)=\sum_aq_ac_1(L_a)$,
whose coordinates are $Kq$, because the first Chern class of an iterated
extension is the sum of those of its quotients. Hence the mass of $q$ is
$(Kq)^T\widetilde G(Kq)$ with $\widetilde G$ positive definite
\cite[Section~2.7]{AGLP12}, and $q$ is massless exactly when $Kq=0$. The
traceless solutions of $Kq=0$ are the multiples of $u$, so a massless
generator exists if and only if $u$ itself is constant on the blocks.
The same conclusion follows from the Higgs mechanism of
\cite{AGL09}: the extension class of a block is a singlet of charge
$e_p-e_q$ with $p,q$ in that block, and its expectation value gives a
mass to every Abelian generator $q$ with $q_p\ne q_q$.
\end{proof}

\begin{lemma}[One side of a wall]\label{lem:wall}
Let $Y$ be a smooth projective threefold, $L,L^\prime$ line bundles on $Y$
and $J_0$ an ample class with $\mu_{J_0}(L)=\mu_{J_0}(L^\prime)$. Let
$0\to L\to E\to L^\prime\to0$ be a non-split extension. There is a
neighbourhood $U$ of $J_0$ in the ample cone such that $E$ is
$\mu_J$-stable for every $J\in U$ with $\mu_J(L)<\mu_J(L^\prime)$.
\end{lemma}
\begin{proof}
A saturated subsheaf $M\subset E$ of rank one is reflexive, hence a line
bundle, as $Y$ is smooth \cite[Section~1.1]{HL10}; it suffices to
show $\mu_J(M)<\mu_J(E)=\frac12(\mu_J(L)+\mu_J(L^\prime))$ for every
such $M$. If the composite $M\to L^\prime$ vanishes, then $M\subset L$,
so $M=L(-D)$ with $D$ effective and $\mu_J(M)\le\mu_J(L)<\mu_J(E)$.
Otherwise $M\to L^\prime$ is injective, $M=L^\prime(-D)$ with $D$
effective, and $D\ne0$, since $D=0$ would split the extension. Put
$\delta_J=\frac12(\mu_J(L^\prime)-\mu_J(L))$, so that
$\mu_J(M)<\mu_J(E)$ means $J^2\cdot D>\delta_J$.

Choose a rational ample class $H$ with $J_0-H$ ample. By Grothendieck's
lemma \cite[Lemma~1.7.9]{HL10} the saturated subsheaves $M\subset E$ with
$\mu_H(M)\ge\mu_H(L^\prime)-1$, that is with $H^2\cdot D\le1$, form a
bounded family, so only finitely many classes $D_1,\dots,D_m\in H^2(Y,\Z)$
occur among them. Each $D_i$ is effective and nonzero, so
$J_0^2\cdot D_i>0$, and $c_1=\min_iJ_0^2\cdot D_i>0$. Let $U$ be a
neighbourhood of $J_0$ in the ample cone on which $J-H$ is ample,
$J^2\cdot D_i>c_1/2$ for all $i$ and $\delta_J<\min(c_1/2,1)$; it exists
by continuity, since $\delta_{J_0}=0$. For $J\in U$ and $M$ in the bounded
family, $J^2\cdot D>c_1/2>\delta_J$. For $M$ outside it, $H^2\cdot D>1$,
and writing $J=H+(J-H)$ with both terms ample gives
$J^2\cdot D\ge H^2\cdot D>1>\delta_J$, because
$H\cdot(J-H)\cdot D$ and $(J-H)^2\cdot D$ are nonnegative on the
effective divisor $D$. In both cases $\mu_J(M)<\mu_J(E)$.
\end{proof}

\begin{proposition}[One block does not suffice]\label{prop:oneblock}
No polystable bundle of slope zero whose shape has exactly one block of
size two, $E=E_1\oplus L_r\oplus L_s\oplus L_o$ with $E_1$ a non-split
extension of $L_q$ by $L_p$, removes the massless generator $u$: either no
such bundle is polystable for any K\"ahler class, or $u_p=u_q$ and $u$
stays massless by Lemma~\ref{lem:branchfactor}.
\end{proposition}
\begin{proof}
Polystability of slope zero requires $\mu_r=\mu_s=\mu_o=0$, that is
$\kappa(J)\perp k_r,k_s,k_o$. If these three rows are linearly
independent, the solutions form the line $\{K^T\kappa=0\}$, since that
line has dimension $4-\rk K=1$ and is contained in them; then
$\mu_p=\mu_q=0$ as well, and $E_1$ has the subbundle $L_p$ of the same
slope, so it is not stable, and the shape does not occur. If the three rows
are dependent, some nonzero $v\in\ker K$ vanishes at $p$ and $q$. Now
$\ker K=\Q\one\oplus\Q u$, so $v=\alpha\one+\beta u$ with
$\alpha+\beta u_p=\alpha+\beta u_q=0$ and $\beta\ne0$, whence $u_p=u_q$.
\end{proof}

Two blocks of size two can do what one cannot. For a shape
$\{p,q\},\{r,s\},\{o\}$ let
\[
 W=\{\mu\in\operatorname{Im}K:\ \mu_o=0,\ \mu_p+\mu_q=0\},
\]
the slopes allowed to a polystable bundle of that shape (the third
condition $\mu_r+\mu_s=0$ follows from $\sum_a\mu_a=0$). When $W$ is a
line $\Q w$, the slopes at every admissible $J$ are $\mu(J)=f(J)w$ for a
scalar $f(J)\ne0$, so the sign of $f$ decides which member of each pair
is the subbundle: the one with $fw_a<0$. We normalise $w$ so that $f>0$
gives the orientation shown in the tables, and we say that a shape is
\emph{realised} at $J$ when the two extension groups
$H^1(L_{\rm sub}L_{\rm quot}^{-1})$ of that orientation are nonzero.

\begin{lemma}[A vanishing on the general tetraquadric]\label{lem:genericvanishing}
For the class $1$ of Table~\ref{tab:eleven}, $h^1(X,L_2L_5^{-1})=
h^1(X,L_5L_2^{-1})=0$ for every tetraquadric $X$ in a nonempty Zariski
open set of smooth members, in particular for the general one.
\end{lemma}
\begin{proof}
Here $k_2-k_5=(-3,-2,2,3)$ and $k_5-k_2=(3,2,-2,-3)$. For $k=(-3,-2,2,3)$
the K\"unneth formula gives $h^1(A,\cO_A(k))=0$ and
$h^2(A,\cO_A(k))=h^2(A,\cO_A(k-2\one))=2\cdot1\cdot3\cdot4=24$, both
groups being $H^1\otimes H^1\otimes H^0\otimes H^0$ in the four factors, so
the restriction sequence gives $h^1(X,\cO_X(k))=\dim\ker\bigl(s\colon
H^2(A,\cO_A(k-2\one))\to H^2(A,\cO_A(k))\bigr)$, where $s$ is the
equation of $X$. In the \v{C}ech description $H^1(\PP^1,\cO(m))$ has the
basis of the monomials $x^ay^b$ with $a,b\le-1$, $a+b=m$, and
$H^0(\PP^1,\cO(m))$ that of the monomials with $a,b\ge0$; a polynomial
acts by multiplication, and a product whose exponents leave the range is
zero. For the quadric
\begin{multline*}
 s_0=x_1^2x_2^2y_3^2x_4^2+y_1^2y_2^2x_3^2x_4^2+y_1^2x_2^2x_3^2y_4^2
 +x_1^2y_2^2y_3^2y_4^2\\+y_1^2(x_2y_2)(x_3y_3)y_4^2+x_1^2(x_2y_2)(x_3y_3)x_4^2
\end{multline*}
each of the $24$ monomials of the source is carried by exactly one of the
six terms to a nonzero monomial of the target, and the $24$ images are
distinct, so the map is a bijection of the two bases; the same holds for
$k=(3,2,-2,-3)$ with the roles of the factors exchanged. Injectivity of a
linear map is an open condition on $s$, and the smooth members form a
nonempty open set, so both hold on a nonempty open set.\footnote{The
program \texttt{code/extension\_branch.py} of the repository verifies the
two permutation matrices and lists, for every datum and every shape with
two blocks of size two, the space $W$, the orientations and the exact
extension dimensions used below; its output is the file
\texttt{extension\_branch\_results.json}.}
\end{proof}

\begin{theorem}[Two blocks]\label{thm:twoblocks}
Let $V$ be one of the eleven data of Table~\ref{tab:eleven} other than the
class $3$.
\begin{enumerate}[label=\textup{(\roman*)},leftmargin=*,itemsep=1pt]
\item For the classes $4$, $5$ and $9$ there is no K\"ahler class $J$ and no
polystable bundle of slope zero of a shape with two blocks of size two on
any smooth tetraquadric. For the class $1$ the same holds on every
tetraquadric of Lemma~\ref{lem:genericvanishing}.
\item For the datum \eqref{eq:bundle} and the classes $2$, $6$, $7$, $8$
and $10$ there are a shape, listed in Table~\ref{tab:twoblocks}, and a
K\"ahler class $J_0$ at which $V$ is polystable, such that the locus
$S=\{J:\mu_o(J)=0,\ \mu_p(J)+\mu_q(J)=0\}$ is a smooth surface near
$J_0$ on which $\mu=fw$ with $f$ changing sign along every curve through
$J_0$ transverse to $\{f=0\}$, and such that for every $J\in S$ with
$f(J)>0$ close enough to $J_0$ and every pair of nonzero classes in the two
extension groups of the table, which are nonzero on every smooth
tetraquadric, the bundle $E_1\oplus E_2\oplus L_o$ is polystable of
slope zero, has the Chern character of $V$, and has structure group
$S(U(2)\times U(2)\times U(1))$ with no massless Abelian gauge boson
beyond $SU(5)$.
\end{enumerate}
\end{theorem}

\begin{table}[t]\small\centering
\caption{The realised shapes of Theorem~\ref{thm:twoblocks}(ii):
the two extensions, the line bundle left alone, the vector $w$ spanning
$W$ (normalised so that $f>0$ makes the stated subbundles of negative
slope), the dimensions of $H^1(L_{\rm sub}L_{\rm quot}^{-1})$ for the two
blocks, the generator $u$ and the coordinates $t$ of $J_0=\sum_it_iJ_i$.
For the class $2$, $r=1.42764\dots$ is the root of
$9r^4-36r^2+14r+16$ in $(1.4,1.5)$, $t_1=(18r^3+27r^2-39r-28)/5=4.7454\dots$
and $t_2=(9r^3+6r^2-22r-4)/10=0.3008\dots$.}\label{tab:twoblocks}
\begin{adjustbox}{max width=\linewidth}
\begin{tabular}{@{}cllccl@{}}
\toprule
datum & blocks & $w$ & $h^1$ & $u$ & $t$\\\midrule
\eqref{eq:bundle} & $L_2\subset E_1\to L_1$, $L_5\subset E_2\to L_3$, $L_4$ & $(1,-1,1,0,-1)$ & $12$, $4$ & $(4,-1,-6,4,-1)$ & $(1,1,1,1)$\\
2 & $L_4\subset E_1\to L_2$, $L_5\subset E_2\to L_3$, $L_1$ & $(0,1,1,-1,-1)$ & $16$, $16$ & $(0,-1,1,0,0)$ & $(t_1,t_2,r,1)$\\
6 & $L_3\subset E_1\to L_2$, $L_5\subset E_2\to L_4$, $L_1$ & $(0,1,-1,1,-1)$ & $16$, $12$ & $(14,-11,-6,4,-1)$ & $(1,1,1,1)$\\
7 & $L_3\subset E_1\to L_1$, $L_5\subset E_2\to L_2$, $L_4$ & $(1,1,-1,0,-1)$ & $4$, $24$ & $(1,-1,0,0,0)$ & $\bigl(1,\tfrac{5+\sqrt{105}}8,\tfrac{7+\sqrt{105}}4,1\bigr)$\\
8 & $L_4\subset E_1\to L_1$, $L_5\subset E_2\to L_2$, $L_3$ & $(1,1,0,-1,-1)$ & $8$, $3$ & $(6,-4,-4,1,1)$ & $(1,1,1,1)$\\
10 & $L_5\subset E_1\to L_1$, $L_4\subset E_2\to L_2$, $L_3$ & $(1,1,0,-1,-1)$ & $12$, $12$ & $(4,-1,-6,4,-1)$ & $(2+\sqrt7,1,2+\sqrt7,1)$\\
\bottomrule
\end{tabular}
\end{adjustbox}
\end{table}

\begin{table}[t]\small\centering
\caption{Theorem~\ref{thm:twoblocks}(i): every shape with two blocks of
size two whose space $W$ is nonzero, for the four excluded data; in all
other shapes $W=0$. In each case $W$ is the line through $w$, and for each
sign of $f$ a group $H^1(L_{\rm sub}L_{\rm quot}^{-1})$ is listed whose
Koszul bound \eqref{eq:koszulbound} is zero, or, marked with an asterisk,
which vanishes by Lemma~\ref{lem:genericvanishing}.}\label{tab:excluded}
\begin{adjustbox}{max width=\linewidth}
\begin{tabular}{@{}clcll@{}}
\toprule
datum & shape & $w$ & vanishing group & vanishing group\\\midrule
1 & $\{L_2,L_4\}$, $\{L_3,L_5\}$, $L_1$ & $(0,1,1,-1,-1)$ & $f>0$: $L_5L_3^{-1}{}^{*}$ & $f<0$: $L_2L_4^{-1}$\\
1 & $\{L_2,L_5\}$, $\{L_3,L_4\}$, $L_1$ & $(0,1,1,-1,-1)$ & $f>0$: $L_5L_2^{-1}{}^{*}$ & $f<0$: $L_3L_4^{-1}$\\
1 & $\{L_1,L_2\}$, $\{L_3,L_5\}$, $L_4$ & $(1,-1,-1,0,1)$ & $f>0$: $L_3L_5^{-1}{}^{*}$ & $f<0$: $L_1L_2^{-1}$\\
1 & $\{L_1,L_3\}$, $\{L_2,L_5\}$, $L_4$ & $(1,-1,-1,0,1)$ & $f>0$: $L_2L_5^{-1}{}^{*}$ & $f<0$: $L_1L_3^{-1}$\\
1 & $\{L_1,L_2\}$, $\{L_3,L_4\}$, $L_5$ & $(-1,1,1,-1,0)$ & $f>0$: $L_1L_2^{-1}$ & $f<0$: $L_3L_4^{-1}$\\
1 & $\{L_1,L_3\}$, $\{L_2,L_4\}$, $L_5$ & $(-1,1,1,-1,0)$ & $f>0$: $L_1L_3^{-1}$ & $f<0$: $L_2L_4^{-1}$\\
4 & $\{L_2,L_4\}$, $\{L_3,L_5\}$, $L_1$ & $(0,-1,-1,1,1)$ & $f>0$: $L_3L_5^{-1}$ & $f<0$: $L_4L_2^{-1}$\\
4 & $\{L_2,L_5\}$, $\{L_3,L_4\}$, $L_1$ & $(0,-1,-1,1,1)$ & $f>0$: $L_2L_5^{-1}$ & $f<0$: $L_4L_3^{-1}$\\
4 & $\{L_1,L_2\}$, $\{L_3,L_5\}$, $L_4$ & $(1,-1,-1,0,1)$ & $f>0$: $L_2L_1^{-1}$ & $f<0$: $L_1L_2^{-1}$\\
4 & $\{L_1,L_3\}$, $\{L_2,L_5\}$, $L_4$ & $(1,-1,-1,0,1)$ & $f>0$: $L_3L_1^{-1}$ & $f<0$: $L_1L_3^{-1}$\\
4 & $\{L_1,L_2\}$, $\{L_3,L_4\}$, $L_5$ & $(1,-1,-1,1,0)$ & $f>0$: $L_2L_1^{-1}$ & $f<0$: $L_1L_2^{-1}$\\
4 & $\{L_1,L_3\}$, $\{L_2,L_4\}$, $L_5$ & $(1,-1,-1,1,0)$ & $f>0$: $L_3L_1^{-1}$ & $f<0$: $L_1L_3^{-1}$\\
5 & $\{L_1,L_4\}$, $\{L_2,L_5\}$, $L_3$ & $(-1,-1,0,1,1)$ & $f>0$: $L_1L_4^{-1}$ & $f<0$: $L_5L_2^{-1}$\\
5 & $\{L_1,L_5\}$, $\{L_2,L_4\}$, $L_3$ & $(-1,-1,0,1,1)$ & $f>0$: $L_1L_5^{-1}$ & $f<0$: $L_5L_1^{-1}$\\
5 & $\{L_1,L_3\}$, $\{L_2,L_5\}$, $L_4$ & $(-1,-1,1,0,1)$ & $f>0$: $L_1L_3^{-1}$ & $f<0$: $L_3L_1^{-1}$\\
5 & $\{L_1,L_5\}$, $\{L_2,L_3\}$, $L_4$ & $(-1,-1,1,0,1)$ & $f>0$: $L_1L_5^{-1}$ & $f<0$: $L_5L_1^{-1}$\\
5 & $\{L_1,L_3\}$, $\{L_2,L_4\}$, $L_5$ & $(1,1,-1,-1,0)$ & $f>0$: $L_3L_1^{-1}$ & $f<0$: $L_1L_3^{-1}$\\
5 & $\{L_1,L_4\}$, $\{L_2,L_3\}$, $L_5$ & $(1,1,-1,-1,0)$ & $f>0$: $L_3L_2^{-1}$ & $f<0$: $L_1L_4^{-1}$\\
9 & $\{L_2,L_3\}$, $\{L_4,L_5\}$, $L_1$ & $(0,1,-1,-1,1)$ & $f>0$: $L_4L_5^{-1}$ & $f<0$: $L_2L_3^{-1}$\\
9 & $\{L_2,L_4\}$, $\{L_3,L_5\}$, $L_1$ & $(0,1,-1,-1,1)$ & $f>0$: $L_3L_5^{-1}$ & $f<0$: $L_2L_4^{-1}$\\
9 & $\{L_1,L_3\}$, $\{L_4,L_5\}$, $L_2$ & $(1,0,-1,-1,1)$ & $f>0$: $L_3L_1^{-1}$ & $f<0$: $L_1L_3^{-1}$\\
9 & $\{L_1,L_4\}$, $\{L_3,L_5\}$, $L_2$ & $(1,0,-1,-1,1)$ & $f>0$: $L_4L_1^{-1}$ & $f<0$: $L_1L_4^{-1}$\\
9 & $\{L_1,L_3\}$, $\{L_2,L_4\}$, $L_5$ & $(1,1,-1,-1,0)$ & $f>0$: $L_3L_1^{-1}$ & $f<0$: $L_1L_3^{-1}$\\
9 & $\{L_1,L_4\}$, $\{L_2,L_3\}$, $L_5$ & $(1,1,-1,-1,0)$ & $f>0$: $L_4L_1^{-1}$ & $f<0$: $L_1L_4^{-1}$\\
\bottomrule
\end{tabular}
\end{adjustbox}
\end{table}

\begin{proof}
(i) A polystable bundle of slope zero and shape $\{p,q\},\{r,s\},\{o\}$
has $\mu(J)\in W$. For each of the four data, Table~\ref{tab:excluded}
lists the shapes with $W\ne0$, and in every one of them $W$ is a line; in
the other shapes all five slopes vanish and the extensions have a
subbundle of equal slope, so they are not stable. On a listed shape
$\mu(J)=fw$ with $f\ne0$, and for each sign of $f$ the subbundle of a
block is its member of negative slope. The table lists for each sign a
group $H^1(L_{\rm sub}L_{\rm quot}^{-1})$ that vanishes, on every smooth
tetraquadric when its bound \eqref{eq:koszulbound} is zero, which is
checked with the K\"unneth formula as in the proof of
Proposition~\ref{prop:eleven}, and on the tetraquadrics of
Lemma~\ref{lem:genericvanishing} in the four cases marked with an
asterisk, which are the groups $H^1(L_2L_5^{-1})$ and $H^1(L_5L_2^{-1})$
of the class $1$ and their copies with $L_3\cong L_2$. The corresponding
block has no non-split extension, so the shape does not occur.

(ii) Fix a row of Table~\ref{tab:twoblocks} and write
$\{p,q\}=\{{\rm sub}_1,{\rm quot}_1\}$, $\{r,s\}$ for the second block.
The slopes vanish at $J_0$: for $t=(1,1,1,1)$ every $\kappa_i$ equals $12$
and every row of $K$ has coordinate sum zero; for the other three $J_0$ the
five polynomials $\mu_a(t)$ vanish at the listed algebraic point, which is
verified by exact arithmetic. The surface $S$ is smooth at $J_0$ when the
differentials of $\mu_o$ and $\mu_p+\mu_q$ are independent there. For
$i\ne j$ one has $\partial\kappa_i/\partial t_j=4(e_1(t)-t_i-t_j)$ and
$\partial\kappa_i/\partial t_i=0$, so at $t=(1,1,1,1)$, using the zero
coordinate sum of the rows, $d\mu_a=-8\sum_jk_a^j\,dt_j$: the two
differentials are $-8$ times the rows $k_o$ and $k_p+k_q$, which are
independent for the three symmetric data (for \eqref{eq:bundle},
$(1,1,-1,-1)$ and $(-2,-3,2,3)$). On $S$ the slopes lie in $W=\Q w$, so
$\mu=fw$ with $f=\mu_{q}/w_q$; its differential along $S$ at $J_0$ is
nonzero, for \eqref{eq:bundle} because $k_1$ is not in the span of $k_4$
and $k_1+k_2$, and for the other data by the same exact computation. Hence
$f$ takes both signs on $S$ arbitrarily close to $J_0$.

For $J\in S$ with $f(J)>0$ the slopes of the subbundles are
$fw_{\rm sub}<0$ and those of the quotients $fw_{\rm quot}>0$, while
$\mu_o=0$ and $\mu_{\rm sub}+\mu_{\rm quot}=0$ on each block. The two
extension groups are nonzero on every smooth tetraquadric, because for each
the bound \eqref{eq:koszulbound} is attained: the K\"unneth groups
$H^1(A,\cO_A(k-2\one))$ and $H^2(A,\cO_A(k))$ vanish, so the restriction
sequence gives $h^1(X,\cO_X(k))=h^1(A,\cO_A(k))+h^2(A,\cO_A(k-2\one))$,
which is the number in the table. Lemma~\ref{lem:wall}, applied to each
block at $J_0$, gives a neighbourhood $U$ of $J_0$ on which both $E_1$ and
$E_2$ are stable for every $J$ with $f(J)>0$; their slopes and that of
$L_o$ are zero, so $E=E_1\oplus E_2\oplus L_o$ is polystable of slope zero
for $J\in S\cap U$, $f(J)>0$. Its Chern character is that of $V$. The
determinants $c_1(E_1)=c_1(L_p)+c_1(L_q)$ and $c_1(E_2)$ differ in every
row, so the two blocks are not isomorphic and the structure group is
$S(U(2)\times U(2)\times U(1))$. Finally $u$ is not constant on the
blocks: in every row $u_{\rm sub}\ne u_{\rm quot}$ for both blocks, as
Table~\ref{tab:twoblocks} shows, so by Lemma~\ref{lem:branchfactor} no
Abelian gauge boson beyond $SU(5)$ is massless at $E$.
\end{proof}

\begin{remark}[What the branch does and does not give]\label{rem:extscope}
Theorem~\ref{thm:twoblocks}(ii) removes the obstruction of
Theorem~\ref{thm:forced} for six of the eleven data, including the fixed
datum, at the cost of leaving the split locus: the polystable K\"ahler
classes of $E$ form a half of the surface $S$, of dimension two, in place
of the ray of Remark~\ref{rem:forcedfixed}, and the extension classes are
further moduli. Everything is stated on the cover $X$. For the quotient
$X/\Gamma$ the K\"ahler class and the two extension classes have to be
$\Gamma$-invariant, which requires the equivariant structures on the line
bundles and the invariant parts of the two groups; for the fixed datum the
group $H^1(L_2L_1^{-1})$ of dimension $12$ has invariant classes by
Section~\ref{sec:pointfailure}, and the invariant part of $H^1(L_5L_3^{-1})$
has not been computed. The charged spectrum and the Yukawa couplings
change along the branch and have not been recomputed; condition C7 is not
addressed, since $S$ and the extension classes are flat directions of the
D-terms; and shapes with a block of size three or more, which by
Lemma~\ref{lem:branchfactor} remove $u$ exactly when $u$ is not constant on
the large block, are not treated. Proposition~\ref{prop:oneblock} and part
(i) close the smallest shapes for the classes $1$, $4$, $5$ and $9$; for
these the Abelian factor can only disappear on a shape with a block of
size at least three.
\end{remark}
